\section{O teste oficial e as resenhas que dividem os modelos}
\label{sec:resultados}
O baseline usa TF-IDF de unigramas e bigramas, regressão logística e texto
completo. A CNN original recebe até 400 tokens. BERT recebe até 192
WordPieces e tem seus pesos ajustados para a tarefa \cite{devlin}.
Cada linha mede o resultado da receita inteira, com seu recorte,
representação e regime de treino.

\begin{tabularx}{\linewidth}{@{}lXrrr@{}}
\toprule
\tabelahead{Divisão} & \tabelahead{Modelo} & \tabelahead{Acurácia} &
\tabelahead{F1 macro} & \tabelahead{CE bruta}\\
\midrule
\LinhasOriginal
\bottomrule
\end{tabularx}
\legendinha{Tabela 5. Escores antes da calibração. Treino com 19.823 resenhas,
validação com 2.478 e teste oficial com 25.000.
\fonte{results/comparison.csv}{Resultados originais}.}

\subsection{Quem acertou quando o outro errou?}
Cruzamos as previsões da CNN e do TF-IDF pelo ID das 25.000 resenhas.
Apenas o TF-IDF acertou \SoTFIDF{} e apenas a CNN acertou \SoCNN{}.
Ambos acertaram 20.821 e ambos erraram 1.589. A fração de acertos exclusivos
do TF-IDF foi 10,87\% no grupo em que a CNN teve recorte, ante 6,29\% no
grupo sem recorte. Para a CNN, as duas frações foram 3,53\% e 3,50\%.
Esse cruzamento compara grupos de textos diferentes. Para ver o efeito da
entrada em um mesmo caso, repetimos as previsões da resenha longa.

\medskip
\painel[SoftBlue]{\rotulo{Resenha longa / 780 tokens / rótulo negativo}\par\smallskip
  \small A crítica explícita aparece depois do limite de 400 tokens. Mantivemos
  os pesos e passamos o prefixo e o texto inteiro aos dois modelos.
  \par\smallskip
  \newcommand{\cutrow}[4]{#1 & #2 & #3\\}
  \begin{tabularx}{\linewidth}{@{}Xrr@{}}
    & \tabelahead{TF-IDF \(p(+)\)} & \tabelahead{CNN \(p(+)\)}\\
    \RecorteVisual
  \end{tabularx}\par\smallskip
  Com limiar de 50\%, a classe dos dois muda de positiva para negativa.
  No texto inteiro, janelas
  como \textit{hate to waste} empurram a margem da CNN para o negativo.}
\smallskip
\painel[SoftTeal]{\rotulo{Resenha curta / 61 tokens / rótulo positivo}\par\smallskip
  \small Ambos leem o texto inteiro. A CNN dá \CurtaCNN{} de escore positivo e
  o TF-IDF, \CurtaTFIDF{}. Na CNN, \textit{part is brilliant} soma \(+0{,}257\)
  à margem e \textit{the ending doesn't make much} soma \(-0{,}167\).
  No TF-IDF, \textit{brilliant} soma \(+0{,}503\) e o bigrama \textit{not one}
  soma \(-0{,}337\), embora esse bigrama esteja numa comparação de enredos.}

\legendinha{Os escores dos dois casos usam as temperaturas originais. As
contribuições acima são para as margens brutas, antes da temperatura. Na CNN,
cada contribuição é o máximo de um filtro vezes o peso desse filtro na diferença
entre os logits. No TF-IDF, é o valor da feature vezes o coeficiente da regressão.
\fonteatual{seminario/diagnosticos.md}{Decomposição e contagens completas}.}

\subsection{Encoders na auditoria sequencial}
Com os mesmos prefixos de até 128 palavras em 2.048 resenhas, fazemos a média
dos estados finais do encoder, excluindo padding e tokens especiais, e
aplicamos uma cabeça linear. Congelados, BERT, DistilBERT e ModernBERT
atingiram 84,52\%, 83,35\% e 81,30\%. Com ajuste completo, chegaram a
88,31\%, 86,80\% e 91,24\%. A CNN e o TF-IDF tiveram 81,95\% e 83,30\%.
\fonte{results/sequences/summary.json}{CE, F1 e configurações por execução}.
